\documentclass[11pt]{article}

\usepackage[margin=1in]{geometry}
\usepackage{amsmath, amssymb}
\usepackage{mathrsfs}
\usepackage{bm}
\usepackage{xcolor}
\usepackage{fancyhdr}
\usepackage{booktabs}
\usepackage{multirow}
\usepackage{array}

\pagestyle{plain}


\begin{document}

% =========== Page 1: test_paper_page1.tex ===========

% ---------- Title ----------
\begin{center}
{\Large\scshape Adam: A Method for Stochastic Optimization}\\[1.5em]
\begin{tabular}{cc}
\textbf{Diederik P. Kingma\textsuperscript{*}} & \textbf{Jimmy Lei Ba\textsuperscript{*}} \\
University of Amsterdam, OpenAI & University of Toronto \\
\texttt{dpkingma@openai.com} & \texttt{jimmy@psi.utoronto.ca}
\end{tabular}
\end{center}

\vspace{1em}

% ---------- Abstract ----------
\begin{center}
{\scshape Abstract}
\end{center}
\begin{quote}
We introduce \emph{Adam}, an algorithm for first-order gradient-based
optimization of stochastic objective functions, based on adaptive
estimates of lower-order moments. The method is straightforward to
implement, is computationally efficient, has little memory
requirements, is invariant to diagonal rescaling of the gradients,
and is well suited for problems that are large in terms of data
and/or parameters. The method is also appropriate for non-stationary
objectives and problems with very noisy and/or sparse gradients. The
hyper-parameters have intuitive interpretations and typically require
little tuning. Some connections to related algorithms, on which
\emph{Adam} was inspired, are discussed. We also analyze the
theoretical convergence properties of the algorithm and provide a
regret bound on the convergence rate that is comparable to the best
known results under the online convex optimization framework.
Empirical results demonstrate that Adam works well in practice and
compares favorably to other stochastic optimization methods. Finally,
we discuss AdaMax, a variant of Adam based on the infinity norm.
\end{quote}

\vspace{1em}

% ---------- Section 1: Introduction ----------
\section*{1 \textsc{Introduction}}

Stochastic gradient-based optimization is of core practical importance
in many fields of science and engineering. Many problems in these
fields can be cast as the optimization of some scalar parameterized
objective function requiring maximization or minimization with respect
to its parameters. If the function is differentiable w.r.t. its
parameters, gradient descent is a relatively efficient optimization
method, since the computation of first-order partial derivatives
w.r.t. all the parameters is of the same computational complexity as
just evaluating the function. Often, objective functions are
stochastic. For example, many objective functions are composed of a
sum of subfunctions evaluated at different subsamples of data; in this
case optimization can be made more efficient by taking gradient steps
w.r.t. individual subfunctions, i.e. stochastic gradient descent (SGD)
or ascent. SGD proved itself as an efficient and effective
optimization method that was central in many machine learning success
stories, such as recent advances in deep learning (Deng et al., 2013;
Krizhevsky et al., 2012; Hinton \& Salakhutdinov, 2006; Hinton et al.,
2012a; Graves et al., 2013). Objectives may also have other sources
of noise than data subsampling, such as dropout (Hinton et al.,
2012b) regularization. For all such noisy objectives, efficient
stochastic optimization techniques are required. The focus of this
paper is on the optimization of stochastic objectives with
high-dimensional parameters spaces. In these cases, higher-order
optimization methods are ill-suited, and discussion in this paper
will be restricted to first-order methods.

We propose \emph{Adam}, a method for efficient stochastic optimization
that only requires first-order gradients with little memory
requirement. The method computes individual adaptive learning rates
for different parameters from estimates of first and second moments
of the gradients; the name \emph{Adam} is derived from adaptive
moment estimation. Our method is designed to combine the advantages
of two recently popular methods: AdaGrad (Duchi et al., 2011), which
works well with sparse gradients, and RMSProp (Tieleman \& Hinton,
2012), which works well in on-line and non-stationary settings;
important connections to these and other stochastic optimization
methods are clarified in section~5. Some of \emph{Adam}'s advantages
are that the magnitudes of parameter updates are invariant to
rescaling of the gradient, its stepsizes are approximately bounded by
the stepsize hyperparameter, it does not require a stationary
objective, it works with sparse gradients, and it naturally performs
a form of step size annealing.

\vspace{2em}

\noindent\rule{0.3\textwidth}{0.4pt}\\[0.3em]
{\footnotesize\textsuperscript{*}Equal contribution. Author ordering
determined by coin flip over a Google Hangout.}

% [page ends here — continues on page 2]

\newpage

% =========== Page 2: test_paper_page2.tex ===========

% ---------- Algorithm 1 box ----------
\noindent
\fbox{\parbox{0.95\textwidth}{\small
\textbf{Algorithm 1:} \emph{Adam}, our proposed algorithm for
stochastic optimization. See section~2 for details, and for a
slightly more efficient (but less clear) order of computation.
$g_t^2$ indicates the elementwise square $g_t \odot g_t$. Good
default settings for the tested machine learning problems are
$\alpha = 0.001$, $\beta_1 = 0.9$, $\beta_2 = 0.999$ and
$\epsilon = 10^{-8}$. All operations on vectors are element-wise.
With $\beta_1^t$ and $\beta_2^t$ we denote $\beta_1$ and $\beta_2$
to the power $t$.\\[0.6em]
\textbf{Require:} $\alpha$: Stepsize\\
\textbf{Require:} $\beta_1, \beta_2 \in [0, 1)$: Exponential decay rates for the moment estimates\\
\textbf{Require:} $f(\theta)$: Stochastic objective function with parameters $\theta$\\
\textbf{Require:} $\theta_0$: Initial parameter vector\\[0.2em]
\hspace*{1em}$m_0 \leftarrow 0$ \hfill (Initialize $1^{\text{st}}$ moment vector)\\
\hspace*{1em}$v_0 \leftarrow 0$ \hfill (Initialize $2^{\text{nd}}$ moment vector)\\
\hspace*{1em}$t \leftarrow 0$ \hfill (Initialize timestep)\\
\hspace*{1em}\textbf{while} $\theta_t$ not converged \textbf{do}\\
\hspace*{2em}$t \leftarrow t + 1$\\
\hspace*{2em}$g_t \leftarrow \nabla_\theta f_t(\theta_{t-1})$ \hfill (Get gradients w.r.t.\ stochastic objective at timestep $t$)\\
\hspace*{2em}$m_t \leftarrow \beta_1 \cdot m_{t-1} + (1 - \beta_1) \cdot g_t$ \hfill (Update biased first moment estimate)\\
\hspace*{2em}$v_t \leftarrow \beta_2 \cdot v_{t-1} + (1 - \beta_2) \cdot g_t^2$ \hfill (Update biased second raw moment estimate)\\
\hspace*{2em}$\widehat{m}_t \leftarrow m_t / (1 - \beta_1^t)$ \hfill (Compute bias-corrected first moment estimate)\\
\hspace*{2em}$\widehat{v}_t \leftarrow v_t / (1 - \beta_2^t)$ \hfill (Compute bias-corrected second raw moment estimate)\\
\hspace*{2em}$\theta_t \leftarrow \theta_{t-1} - \alpha \cdot \widehat{m}_t / (\sqrt{\widehat{v}_t} + \epsilon)$ \hfill (Update parameters)\\
\hspace*{1em}\textbf{end while}\\
\hspace*{1em}\textbf{return} $\theta_t$ (Resulting parameters)
}}

\vspace{1em}

% ---------- Body prose ----------
\noindent
In section~2 we describe the algorithm and the properties of its
update rule. Section~3 explains our initialization bias correction
technique, and section~4 provides a theoretical analysis of Adam's
convergence in online convex programming. Empirically, our method
consistently outperforms other methods for a variety of models and
datasets, as shown in section~6. Overall, we show that Adam is a
versatile algorithm that scales to large-scale high-dimensional
machine learning problems.

\section*{2 \textsc{Algorithm}}

See algorithm~1 for pseudo-code of our proposed algorithm
\emph{Adam}. Let $f(\theta)$ be a noisy objective function: a
stochastic scalar function that is differentiable w.r.t. parameters
$\theta$. We are interested in minimizing the expected value of
this function, $\mathbb{E}[f(\theta)]$ w.r.t. its parameters
$\theta$. With $f_1(\theta), \ldots, f_T(\theta)$ we denote the
realisations of the stochastic function at subsequent timesteps
$1, \ldots, T$. The stochasticity might come from the evaluation
at random subsamples (minibatches) of datapoints, or arise from
inherent function noise. With
$g_t = \nabla_\theta f_t(\theta)$ we denote the gradient, i.e.
the vector of partial derivatives of $f_t$, w.r.t $\theta$
evaluated at timestep $t$.

The algorithm updates exponential moving averages of the gradient
($m_t$) and the squared gradient ($v_t$) where the hyper-parameters
$\beta_1, \beta_2 \in [0, 1)$ control the exponential decay rates of
these moving averages. The moving averages themselves are estimates
of the $1^\text{st}$ moment (the mean) and the $2^\text{nd}$ raw
moment (the uncentered variance) of the gradient. However, these
moving averages are initialized as (vectors of) 0's, leading to
moment estimates that are biased towards zero, especially during
the initial timesteps, and especially when the decay rates are
small (i.e.\ the $\beta$s are close to 1). The good news is that
this initialization bias can be easily counteracted, resulting in
bias-corrected estimates $\widehat{m}_t$ and $\widehat{v}_t$. See
section~3 for more details.

Note that the efficiency of algorithm~1 can, at the expense of
clarity, be improved upon by changing the order of computation,
e.g.\ by replacing the last three lines in the loop with the
following lines:
$\alpha_t = \alpha \cdot \sqrt{1 - \beta_2^t}/(1 - \beta_1^t)$
and
$\theta_t \leftarrow \theta_{t-1} - \alpha_t \cdot m_t/(\sqrt{v_t} + \hat{\epsilon})$.

\subsection*{2.1 \textsc{Adam's update rule}}

An important property of Adam's update rule is its careful choice
of stepsizes. Assuming $\epsilon = 0$, the effective step taken in
parameter space at timestep $t$ is
$\Delta_t = \alpha \cdot \widehat{m}_t / \sqrt{\widehat{v}_t}$.
The effective stepsize has two upper bounds:
$|\Delta_t| \leq \alpha \cdot (1 - \beta_1) / \sqrt{1 - \beta_2}$ in
the case $(1 - \beta_1) > \sqrt{1 - \beta_2}$, and
$|\Delta_t| \leq \alpha$

% [page ends here — continues on page 3]

\newpage

% =========== Page 3: test_paper_page3.tex ===========
\setcounter{equation}{0}

% ---------- Continuation of Sec 2.1 ----------
\noindent
otherwise. The first case only happens in the most severe case of
sparsity: when a gradient has been zero at all timesteps except at
the current timestep. For less sparse cases, the effective stepsize
will be smaller. When
$(1 - \beta_1) = \sqrt{1 - \beta_2}$ we have that
$|\widehat{m}_t / \sqrt{\widehat{v}_t}| < 1$ therefore
$|\Delta_t| < \alpha$. In more common scenarios, we will have that
$\widehat{m}_t / \sqrt{\widehat{v}_t} \approx \pm 1$ since
$|\mathbb{E}[g] / \sqrt{\mathbb{E}[g^2]}| \leq 1$. The effective
magnitude of the steps taken in parameter space at each timestep
are approximately bounded by the stepsize setting $\alpha$, i.e.,
$|\Delta_t| \lessapprox \alpha$. This can be understood as
establishing a \emph{trust region} around the current parameter
value, beyond which the current gradient estimate does not provide
sufficient information. This typically makes it relatively easy to
know the right scale of $\alpha$ in advance. For many machine
learning models, for instance, we often know in advance that good
optima with high probability within some set region in parameter
space; it is not uncommon, for example, to have a prior
distribution over the parameters. Since $\alpha$ sets (an upper
bound of) the magnitude of steps in parameter space, we can often
deduce the right order of magnitude of $\alpha$ such that optima
can be reached from $\theta_0$ within some number of iterations.
It may be that with a slight abuse of terminology, we will call
the ratio $\widehat{m}_t / \sqrt{\widehat{v}_t}$ the
\emph{signal-to-noise} ratio ($SNR$). With a smaller SNR the
effective stepsize $\Delta_t$ will be closer to zero. This is a
desirable property, since a smaller SNR means that there is
greater uncertainty about whether the direction of
$\widehat{m}_t$ corresponds to the direction of the true gradient.
For example, the SNR value typically becomes closer to 0 towards
an optimum, leading to smaller effective steps in parameter space:
a form of automatic annealing. The effective stepsize $\Delta_t$
is also invariant to the scale of the gradients; rescaling the
gradients $g$ with factor $c$ will scale $\widehat{m}_t$ with a
factor $c$ and $\widehat{v}_t$ with a factor $c^2$, which cancel
out:
$(c \cdot \widehat{m}_t) / (\sqrt{c^2 \cdot \widehat{v}_t}) =
 \widehat{m}_t / \sqrt{\widehat{v}_t}$.

\vspace{1em}

% ---------- Section 3: Initialization Bias Correction ----------
\section*{3 \textsc{Initialization bias correction}}

As explained in section~2, Adam utilizes initialization bias
correction terms. We will here derive the term for the second
moment estimate; the derivation for the first moment estimate is
completely analogous. Let $g$ be the gradient of the stochastic
objective $f$, and we wish to estimate its second raw moment
(uncentered variance) using an exponential moving average of the
squared gradient, with decay rate $\beta_2$. Let
$g_1, \ldots, g_T$ be the gradients at subsequent timesteps, each
a draw from an underlying gradient distribution
$g_t \sim p(g_t)$. Let us initialize the exponential moving
average as $v_0 = 0$ (a vector of zeros). First note that the
update at timestep $t$ of the exponential moving average
$v_t = \beta_2 \cdot v_{t-1} + (1 - \beta_2) \cdot g_t^2$ (where
$g_t^2$ indicates the elementwise square $g_t \odot g_t$) can be
written as a function of the gradients at all previous timesteps:
%
\begin{equation}
v_t = (1 - \beta_2) \sum_{i=1}^{t} \beta_2^{t-i} \cdot g_i^2
\end{equation}
%
We wish to know how $\mathbb{E}[v_t]$, the expected value of the
exponential moving average at timestep $t$, relates to the true
second moment $\mathbb{E}[g_t^2]$, so we can correct for the
discrepancy between the two. Taking expectations of the left-hand
and right-hand sides of eq.~(1):
%
\begin{align}
\mathbb{E}[v_t] &= \mathbb{E}\!\left[(1 - \beta_2) \sum_{i=1}^{t} \beta_2^{t-i} \cdot g_i^2\right] \\
                &= \mathbb{E}[g_t^2] \cdot (1 - \beta_2) \sum_{i=1}^{t} \beta_2^{t-i} + \zeta \\
                &= \mathbb{E}[g_t^2] \cdot (1 - \beta_2^t) + \zeta
\end{align}
%
where $\zeta = 0$ if the true second moment $\mathbb{E}[g_t^2]$ is
stationary; otherwise $\zeta$ can be kept small since the
exponential decay rate $\beta_1$ can (and should) be chosen such
that the exponential moving average assigns small weights to
gradients too far in the past. What is left is the term
$(1 - \beta_2^t)$ which is caused by initializing the running
average with zeros. In algorithm~1 we therefore divide by this
term to correct the initialization bias.

In case of sparse gradients, for a reliable estimate of the second
moment one needs to average over many gradients by choosing a
small value of $\beta_2$; however it is exactly this case of small
$\beta_2$ where a lack of initialisation bias correction would
lead to initial steps that are much larger.

% [page ends here — continues on page 4]

\newpage

% =========== Page 4: test_paper_page4.tex ===========
\setcounter{equation}{4}

\section*{4 \textsc{Convergence analysis}}

We analyze the convergence of Adam using the online learning
framework proposed in (Zinkevich, 2003). Given an arbitrary,
unknown sequence of convex cost functions
$f_1(\theta), f_2(\theta), \ldots, f_T(\theta)$. At each time $t$,
our goal is to predict the parameter $\theta_t$ and evaluate it on
a previously unknown cost function $f_t$. Since the nature of the
sequence is unknown in advance, we evaluate our algorithm using
the \emph{regret}, that is the sum of all the previous difference
between the online prediction $f_t(\theta_t)$ and the best fixed
point parameter $f_t(\theta^*)$ from a feasible set $\mathcal{X}$
for all the previous steps. Concretely, the regret is defined as:
%
\begin{equation}
R(T) = \sum_{t=1}^{T} [f_t(\theta_t) - f_t(\theta^*)]
\end{equation}
%
where $\theta^* = \arg\min_{\theta \in \mathcal{X}} \sum_{t=1}^{T} f_t(\theta)$.
We show Adam has $O(\sqrt{T})$ regret bound and a proof is given
in the appendix. Our result is comparable to the best known bound
for this general convex online learning problem. We also use some
definitions simplify our notation, where
$g_t \triangleq \nabla f_t(\theta_t)$ and $g_{t,i}$ as the
$i^\text{th}$ element. We define
$g_{1:t,i} \in \mathbb{R}^t$ as a vector that contains the
$i^\text{th}$ dimension of the gradients over all iterations till
$t$, $g_{1:t,i} = [g_{1,i}, g_{2,i}, \ldots, g_{t,i}]$. Also, we
define $\gamma \triangleq \frac{\beta_1^2}{\sqrt{\beta_2}}$. Our
following theorem holds when the learning rate $\alpha_t$ is
decaying at a rate of $t^{-\frac{1}{2}}$ and first moment running
average coefficient $\beta_{1,t}$ decay exponentially with
$\lambda$, that is typically close to 1, e.g.\
$1 - 10^{-8}$.

\vspace{0.5em}

\noindent\textbf{Theorem 4.1.} \emph{Assume that the function
$f_t$ has bounded gradients,
$\|\nabla f_t(\theta)\|_2 \leq G$,
$\|\nabla f_t(\theta)\|_\infty \leq G_\infty$ for all
$\theta \in \mathbb{R}^d$ and distance between any $\theta_t$
generated by Adam is bounded,
$\|\theta_n - \theta_m\|_2 \leq D$,
$\|\theta_m - \theta_n\|_\infty \leq D_\infty$ for any
$m, n \in \{1, \ldots, T\}$, and $\beta_1, \beta_2 \in [0, 1)$
satisfy $\frac{\beta_1^2}{\sqrt{\beta_2}} < 1$. Let
$\alpha_t = \frac{\alpha}{\sqrt{t}}$ and
$\beta_{1,t} = \beta_1 \lambda^{t-1}$, $\lambda \in (0, 1)$.
Adam achieves the following guarantee, for all $T \geq 1$.}
%
\[
R(T) \leq \frac{D^2}{2\alpha(1 - \beta_1)}
          \sum_{i=1}^{d} \sqrt{T \, \widehat{v}_{T,i}}
        + \frac{\alpha(1 + \beta_1) G_\infty}
               {(1 - \beta_1)\sqrt{1 - \beta_2}\,(1 - \gamma)^2}
          \sum_{i=1}^{d} \|g_{1:T,i}\|_2
        + \sum_{i=1}^{d}
          \frac{D_\infty^2 \, G_\infty \sqrt{1 - \beta_2}}
               {2\alpha(1 - \beta_1)(1 - \lambda)^2}
\]

\noindent
Our Theorem~4.1 implies when the data features are sparse and
bounded gradients, the summation term can be much smaller than
its upper bound
$\sum_{i=1}^{d} \|g_{1:T,i}\|_2 \ll dG_\infty\sqrt{T}$ and
$\sum_{i=1}^{d} \sqrt{T\,\widehat{v}_{T,i}} \ll dG_\infty\sqrt{T}$,
in particular if the class of function and data features are in
the form of section~1.2 in (Duchi et al., 2011). Their results
for the expected value
$\mathbb{E}[\sum_{i=1}^{d} \|g_{1:T,i}\|_2]$ also apply to Adam.
In particular, the adaptive method, such as Adam and Adagrad, can
achieve $O(\log d\sqrt{T})$, an improvement over $O(\sqrt{dT})$
for the non-adaptive method. Decoupling $\beta_1$ towards zero is
important in our theoretical analysis and also matches previous
empirical findings, e.g.\ (Sutskever et al., 2013) suggests
reducing the momentum coefficient in the end of training can
improve convergence.

Finally, we can show the average regret of Adam converges,

\vspace{0.5em}

\noindent\textbf{Corollary 4.2.} \emph{Assume that the function
$f_t$ has bounded gradients,
$\|\nabla f_t(\theta)\|_2 \leq G$,
$\|\nabla f_t(\theta)\|_\infty \leq G_\infty$ for all
$\theta \in \mathbb{R}^d$ and distance between any $\theta_t$
generated by Adam is bounded,
$\|\theta_n - \theta_m\|_2 \leq D$,
$\|\theta_m - \theta_n\|_\infty \leq D_\infty$ for any
$m, n \in \{1, \ldots, T\}$. Adam achieves the following
guarantee, for all $T \geq 1$.}
%
\[
\frac{R(T)}{T} = O\!\left(\frac{1}{\sqrt{T}}\right)
\]

\noindent
This result can be obtained by using Theorem~4.1 and
$\sum_{i=1}^{d} \|g_{1:T,i}\|_2 \leq dG_\infty\sqrt{T}$. Thus,
$\lim_{T \to \infty} \frac{R(T)}{T} = 0$.

\vspace{0.5em}

\section*{5 \textsc{Related work}}

Optimization methods having a direct relation to Adam are RMSProp
(Tieleman \& Hinton, 2012; Graves, 2013) and AdaGrad (Duchi et al.,
2011); these relationships are discussed below. Other stochastic
optimization methods include vSGD (Schaul et al., 2012), AdaDelta
(Zeiler, 2012) and the natural Newton method from Roux \&
Fitzgibbon (2010), all setting stepsizes by estimating curvature

% [page ends here — continues on page 5]

\newpage

% =========== Page 5: test_paper_page5.tex ===========
\setcounter{equation}{5}

\noindent
from first-order information. The Sum-of-Functions Optimizer
(SFO) (Sohl-Dickstein et al., 2014) is a quasi-Newton method
based on minibatches, but (unlike Adam) has memory requirements
linear in the number of minibatch partitions of a dataset, which
is often infeasible on memory-constrained systems such as a GPU.
Like natural gradient descent (NGD) (Amari, 1998), Adam employs
a preconditioner that adapts to the geometry of the data, since
$\widehat{v}_t$ is an approximation to the diagonal of the Fisher
information matrix (Pascanu \& Bengio, 2013); however, Adam's
preconditioner (like AdaGrad's) is more conservative in its
adaption than vanilla NGD by preconditioning with the square root
of the inverse of the diagonal Fisher information matrix
approximation.

\vspace{0.5em}

\noindent\textbf{RMSProp:} An optimization method closely related
to Adam is RMSProp (Tieleman \& Hinton, 2012). A version with
momentum has sometimes been used (Graves, 2013). There are a few
important differences between RMSProp with momentum and Adam:
RMSProp with momentum generates its parameter updates using a
momentum on the rescaled gradient, whereas Adam updates are
directly estimated using a running average of first and second
moment of the gradient. RMSProp also lacks a bias-correction
term; this matters most in case of a value of $\beta_2$ close to
$1$ (required in case of sparse gradients), since in that case
not correcting the bias leads to very large stepsizes and often
divergence, as we also empirically demonstrate in section~6.4.

\vspace{0.5em}

\noindent\textbf{AdaGrad:} An algorithm that works well for
sparse gradients is AdaGrad (Duchi et al., 2011). Its basic
version updates parameters as
$\theta_{t+1} = \theta_t - \alpha \cdot g_t / \sqrt{\sum_{i=1}^{t} g_i^2}$.
Note that if we choose $\beta_2$ to be infinitesimally close to
$1$ from below, then
$\lim_{\beta_2 \to 1} \widehat{v}_t = t^{-1} \cdot \sum_{i=1}^{t} g_i^2$.
AdaGrad corresponds to a version of Adam with $\beta_1 = 0$,
infinitesimal $(1 - \beta_2)$ and a replacement of $\alpha$ by an
annealed version $\alpha_t = \alpha \cdot t^{-1/2}$, namely
$\theta_t - \alpha \cdot t^{-1/2} \cdot \widehat{m}_t / \sqrt{\lim_{\beta_2 \to 1} \widehat{v}_t}
 = \theta_t - \alpha \cdot t^{-1/2} \cdot g_t / \sqrt{t^{-1} \cdot \sum_{i=1}^{t} g_i^2}
 = \theta_t - \alpha \cdot g_t / \sqrt{\sum_{i=1}^{t} g_i^2}$.
Note that this direct correspondence between Adam and AdaGrad
does not hold when removing the bias-correction terms; without
bias correction, like in RMSProp, a $\beta_2$ infinitesimally
close to $1$ would lead to infinitely large bias, and infinitely
large parameter updates.

\vspace{0.5em}

\section*{6 \textsc{Experiments}}

To empirically evaluate the proposed method, we investigated
different popular machine learning models, including logistic
regression, multilayer fully connected neural networks and deep
convolutional neural networks. Using large models and datasets,
we demonstrate Adam can efficiently solve practical deep learning
problems.

We use the same parameter initialization when comparing different
optimization algorithms. The hyper-parameters, such as learning
rate and momentum, are searched over a dense grid and the results
are reported using the best hyper-parameter setting.

\subsection*{6.1 \textsc{Experiment: Logistic regression}}

We evaluate our proposed method on L2-regularized multi-class
logistic regression using the MNIST dataset. Logistic regression
has a well-studied convex objective, making it suitable for
comparison of different optimizers without worrying about local
minimum issues. The stepsize $\alpha$ in our logistic regression
experiments is adjusted by $1/\sqrt{t}$ decay, namely
$\alpha_t = \frac{\alpha}{\sqrt{t}}$ that matches with our
theoratical prediction from section~4. The logistic regression
classifies the class label directly on the 784 dimension image
vectors. We compare Adam to accelerated SGD with Nesterov
momentum and Adagrad using minibatch size of 128. According to
Figure~1, we found that the Adam yields similar convergence as
SGD with momentum and both converge faster than Adagrad.

As discussed in (Duchi et al., 2011), Adagrad can efficiently
deal with sparse features and gradients as one of its main
theoretical results whereas SGD is low at learning rare features.
Adam with $1/\sqrt{t}$ decay on its stepsize should theoretically
match the performance of Adagrad. We examine the sparse feature
problem using IMDB movie review dataset from (Maas et al., 2011).
We pre-process the IMDB movie reviews into bag-of-words (BoW)
feature vectors including the first 10,000 most frequent words.
The 10,000 dimension BoW feature vector for each review is highly
sparse. As suggested in (Wang \& Manning, 2013), 50\% dropout
noise can be applied to the BoW features during

% [page ends here — continues on page 6]

\newpage

% =========== Page 6: test_paper_page6.tex ===========
\setcounter{equation}{5}

% ---------- Figure 1 placeholder ----------
\noindent
\begin{center}
\fbox{\parbox{0.92\textwidth}{\centering\vspace{1em}
\textbf{[Figure 1 --- placeholder]}\\[0.8em]
Two side-by-side training-cost curves for L2-regularized logistic
regression, showing Adam vs.\ baseline optimizers.\\[0.4em]
\textbf{Left panel --- MNIST Logistic Regression}: training cost
(y-axis, $\approx 0.2$ to $0.7$) vs.\ iterations over entire
dataset (x-axis, $0$ to $\sim 45$). Three curves: \emph{AdaGrad}
(slowest), \emph{SGDNesterov}, and \emph{Adam} (both SGDNesterov
and Adam converge faster than AdaGrad; Adam and SGDNesterov are
visually comparable).\\[0.3em]
\textbf{Right panel --- IMDB BoW feature Logistic Regression}:
training cost (y-axis, $\approx 0.25$ to $0.50$) vs.\ iterations
over entire dataset (x-axis, $0$ to $\sim 160$). Four curves, all
with dropout applied to the sparse BoW features: \emph{Adagrad +
dropout}, \emph{RMSProp + dropout}, \emph{SGDNesterov + dropout},
and \emph{Adam + dropout}. The RMSProp curve is visibly noisier
(larger fluctuations) than the others; Adam + dropout converges
to the lowest training cost overall, with Adagrad+dropout
tracking closely.
\vspace{1em}}}
\end{center}
\vspace{0.3em}
\begin{center}
\textbf{Figure 1:} Logistic regression training negative log
likelihood on MNIST images and IMDB movie reviews with 10,000
bag-of-words (BoW) feature vectors.
\end{center}

\vspace{1em}

% ---------- Prose after Fig. 1 ----------
\noindent
training to prevent over-fitting. In figure~1, Adagrad outperforms
SGD with Nesterov momentum by a large margin both with and
without dropout noise. Adam converges as fast as Adagrad. The
empirical performance of Adam is consistent with our theoretical
findings in sections~2 and~4. Similar to Adagrad, Adam can take
advantage of sparse features and obtain faster convergence rate
than normal SGD with momentum.

\vspace{0.5em}

\subsection*{6.2 \textsc{Experiment: Multi-layer neural networks}}

Multi-layer neural network are powerful models with non-convex
objective functions. Although our convergence analysis does not
apply to non-convex problems, we empirically found that Adam
often outperforms other methods in such cases. In our
experiments, we made model choices that are consistent with
previous publications in the area; a neural network model with
two fully connected hidden layers with $1000$ hidden units each
and ReLU activation are used for this experiment with minibatch
size of $128$.

First, we study different optimizers using the standard
deterministic cross-entropy objective function with $L_2$ weight
decay on the parameters to prevent over-fitting. The
sum-of-functions (SFO) method (Sohl-Dickstein et al., 2014) is a
recently proposed quasi-Newton method that works with minibatches
of data and has shown good performance on optimization of
multi-layer neural networks. We used their implementation and
compared with Adam to train such models. Figure~2 shows that
Adam makes faster progress in terms of both the number of
iterations and wall-clock time. Due to the cost of updating
curvature information, SFO is 5--10x slower per iteration
compared to Adam, and has a memory requirement that is linear in
the number minibatches.

Stochastic regularization methods, such as dropout, are an
effective way to prevent over-fitting and often used in practice
due to their simplicity. SFO assumes deterministic
subfunctions, and indeed failed to converge on cost functions
with stochastic regularization. We compare the effectiveness of
Adam to other stochastic first order methods on multi-layer
neural networks trained with dropout noise. Figure~2 shows our
results; Adam shows better convergence than other methods.

\vspace{0.5em}

\subsection*{6.3 \textsc{Experiment: Convolutional neural networks}}

Convolutional neural networks (CNNs) with several layers of
convolution, pooling and non-linear units have shown considerable
success in computer vision tasks. Unlike most fully connected
neural nets, weight sharing in CNNs results in vastly different
gradients in different layers. A smaller learning rate for the
convolution layers is often used in practice when applying SGD.
We show the effectiveness of Adam in deep CNNs. Our CNN
architecture has three alternating stages of $5 \times 5$
convolution filters and $3 \times 3$ max pooling with stride of
$2$ that are followed by a fully connected layer of $1000$
rectified linear hidden units (ReLU's). The input image are
pre-processed by whitening, and

% [page ends here — continues on page 7]

\newpage

% =========== Page 7: test_paper_page7.tex ===========
\setcounter{equation}{5}

% ---------- Figure 2 placeholder ----------
\noindent
\begin{center}
\fbox{\parbox{0.92\textwidth}{\centering\vspace{1em}
\textbf{[Figure 2 --- placeholder]}\\[0.8em]
Two-panel comparison of optimizers on MNIST multilayer neural
networks.\\[0.4em]
\textbf{(a) MNIST Multilayer Neural Network + dropout}: training
cost on log-scale y-axis ($\sim 10^{-1}$ to $10^1$) vs.\ iterations
over entire dataset (0 to $\sim$200). Five curves: AdaGrad,
RMSProp, SGDNesterov, AdaDelta, Adam. Adam reaches the lowest
training cost, visible as the bottom curve after
$\sim$50 iterations.\\[0.4em]
\textbf{(b) MNIST Multilayer Neural Network (deterministic)}:
two sub-sub-panels comparing the sum-of-functions (SFO) optimizer
to Adam. The top sub-panel plots training error vs.\ iterations;
the bottom sub-panel plots training error vs.\ \emph{normalized
epochs} (wall-clock / curvature-update-cost). Adam outperforms
SFO on both axes --- faster per iteration \emph{and} faster in
wall-clock time, because SFO pays 5--10x more per iteration for
its curvature update and has memory requirements linear in the
number of minibatch partitions.
\vspace{1em}}}
\end{center}
\vspace{0.3em}
\begin{center}
\textbf{Figure 2:} Training of multilayer neural networks on
MNIST images. (a) Neural networks using dropout stochastic
regularization. (b) Neural networks with deterministic cost
function. We compare with the sum-of-functions (SFO) optimizer
(Sohl-Dickstein et al., 2014).
\end{center}

\vspace{1em}

% ---------- Figure 3 placeholder ----------
\noindent
\begin{center}
\fbox{\parbox{0.92\textwidth}{\centering\vspace{1em}
\textbf{[Figure 3 --- placeholder]}\\[0.8em]
Two-panel CIFAR-10 ConvNet comparison, architecture
\texttt{c64-c64-c128-1000}.\\[0.4em]
\textbf{(left) First 3 epochs}: training cost (linear y-axis,
$\sim 0.5$ to $3.0$) vs.\ iterations (0 to 3 epochs). Six
curves: AdaGrad, AdaGrad+dropout, SGDNesterov,
SGDNesterov+dropout, Adam, Adam+dropout. In the early phase,
AdaGrad and Adam both make rapid initial progress and are the
fastest.\\[0.4em]
\textbf{(right) All 45 epochs}: training cost on log y-axis
($10^{-4}$ to $10^1$) vs.\ iterations over entire dataset (0 to
45 epochs). Same six curves. \textbf{Adam+dropout reaches the
lowest training cost} by a visible margin over SGDNesterov
variants. AdaGrad's early lead \emph{reverses}: by epoch 15 it
is visibly the slowest, while SGDNesterov and Adam pull ahead.
This is the sign-change the surrounding prose explains: for
CNNs, Adam's second-moment preconditioning stops helping after
a few epochs.
\vspace{1em}}}
\end{center}
\vspace{0.3em}
\begin{center}
\textbf{Figure 3:} Convolutional neural networks training cost.
(left) Training cost for the first three epochs. (right)
Training cost over 45 epochs. CIFAR-10 with c64-c64-c128-1000
architecture.
\end{center}

\vspace{1em}

% ---------- Body prose (continuation from Sec 6.3) ----------
\noindent
dropout noise is applied to the input layer and fully connected
layer. The minibatch size is also set to $128$ similar to
previous experiments.

Interestingly, although both Adam and Adagrad make rapid
progress lowering the cost in the initial stage of the training,
shown in Figure~3 (left), Adam and SGD eventually converge
considerably faster than Adagrad for CNNs shown in Figure~3
(right). We notice the second moment estimate $\widehat{v}_t$
vanishes to zeros after a few epochs and is dominated by the
$\epsilon$ in algorithm~1. The second moment estimate is
therefore a poor approximation to the geometry of the cost
function in CNNs comparing to fully connected network from
Section~6.2. Whereas, reducing the minibatch variance through
the first moment is more important in CNNs and contributes to
the speed-up. As a result, Adagrad converges much slower than
others in this particular experiment. Though Adam shows marginal
improvement over SGD with momentum, it adapts learning rate
scale for different layers instead of hand picking manually as
in SGD.

% [page ends here — continues on page 8]

\newpage

% =========== Page 8: test_paper_page8.tex ===========
\setcounter{equation}{5}

% ---------- Figure 4 placeholder ----------
\noindent
\begin{center}
\fbox{\parbox{0.92\textwidth}{\centering\vspace{1em}
\textbf{[Figure 4 --- placeholder]}\\[0.8em]
\textbf{2$\times$6 grid} of VAE training-loss curves for a
parameter sweep over $\beta_1 \in \{0, 0.9\}$,
$\beta_2 \in \{0.99, 0.999, 0.9999\}$, and
$\log_{10}(\alpha) \in [-5, -1]$, comparing Adam \emph{with}
bias-correction terms (\textbf{red line}) vs.\ Adam
\emph{without} bias-correction terms (\textbf{green line}).
Y-axes plot loss ($\sim 100$ to $140$); x-axes plot
$\log_{10}(\alpha)$.\\[0.4em]
\textbf{(a) after 10 epochs (left three columns)}: at
$\beta_2 = 0.99$ the two curves are nearly coincident; at
$\beta_2 = 0.999$ and especially $\beta_2 = 0.9999$, the green
(no-correction) curve lies \emph{above} the red (with-correction)
curve --- bias correction reduces loss. The gap grows as
$\beta_2 \to 1$.\\[0.4em]
\textbf{(b) after 100 epochs (right three columns)}: same pattern
but the gap is larger and more systematic. Without bias
correction and with $\beta_2$ close to $1$, training is visibly
worse --- the green curve rises at small $\alpha$ values where
the red curve stays flat, indicating the no-bias-correction
version is unstable in the regime that's actually needed for
sparse gradients.\\[0.4em]
\textbf{Takeaway}: bias correction matters most exactly when
$\beta_2 \to 1$, which is the regime required for sparse
gradients. Removing the correction turns Adam into a version of
momentum-RMSProp that is numerically unstable in this useful
corner of parameter space.
\vspace{1em}}}
\end{center}
\vspace{0.3em}
\begin{center}
\textbf{Figure 4:} Effect of bias-correction terms (red line)
versus no bias correction terms (green line) after 10 epochs
(left) and 100 epochs (right) on the loss (y-axes) when learning
a Variational Auto-Encoder (VAE) (Kingma \& Welling, 2013), for
different settings of stepsize $\alpha$ (x-axes) and
hyper-parameters $\beta_1$ and $\beta_2$.
\end{center}

\vspace{1em}

% ---------- Section 6.4 ----------
\subsection*{6.4 \textsc{Experiment: Bias-correction term}}

We also empirically evaluate the effect of the bias correction
terms explained in sections~2 and~3. Discussed in section~5,
removal of the bias correction terms results in a version of
RMSProp (Tieleman \& Hinton, 2012) with momentum. We vary the
$\beta_1$ and $\beta_2$ when training a variational auto-encoder
(VAE) (Kingma \& Welling, 2013) with the same architecture as in
(Kingma \& Welling, 2013) with a single hidden layer with 500
hidden units with softplus nonlinearities and a 50-dimensional
spherical Gaussian latent variable. We iterated over a broad
range of hyper-parameter choices, i.e.\
$\beta_1 \in [0, 0.9]$ and
$\beta_2 \in [0.99, 0.999, 0.9999]$, and
$\log_{10}(\alpha) \in [-5, \ldots, -1]$. Values of $\beta_2$
close to $1$, required for robustness to sparse gradients,
results in larger initialization bias; therefore we expect the
bias correction term is important in such cases of slow decay,
preventing an adverse effect on optimization.

In Figure~4, values $\beta_2$ close to $1$ indeed lead to
instabilities in training when no bias correction term was
present, especially at first few epochs of the training. The
best results were achieved with small values of $(1 - \beta_2)$
and bias correction; this was more apparent towards the end of
optimization when gradients tends to become sparser as hidden
units specialize to specific patterns. In summary, Adam performed
equal or better than RMSProp, regardless of hyper-parameter
setting.

\vspace{0.5em}

\section*{7 \textsc{Extensions}}

\subsection*{7.1 \textsc{AdaMax}}

In Adam, the update rule for individual weights is to scale
their gradients inversely proportional to a (scaled) $L^2$ norm
of their individual current and past gradients. We can
generalize the $L^2$ norm based update rule to a $L^p$ norm
based update rule. Such variants become numerically unstable for
large $p$. However, in the special case where we let
$p \to \infty$, a surprisingly simple and stable algorithm
emerges; see algorithm~2. We'll now derive the algorithm. Let,
in case of the $L^p$ norm, the stepsize at time $t$ be inversely
proportional to $v_t^{1/p}$, where:
%
\begin{align}
v_t &= \beta_2^p \, v_{t-1} + (1 - \beta_2^p) |g_t|^p \\
    &= (1 - \beta_2^p) \sum_{i=1}^{t} \beta_2^{p(t-i)} \cdot |g_i|^p
\end{align}

% [page ends here — continues on page 9]

\newpage

% =========== Page 9: test_paper_page9.tex ===========
\setcounter{equation}{7}

% ---------- Algorithm 2 box ----------
\noindent
\fbox{\parbox{0.95\textwidth}{\small
\textbf{Algorithm 2:} \emph{AdaMax}, a variant of Adam based on
the infinity norm. See section~7.1 for details. Good default
settings for the tested machine learning problems are
$\alpha = 0.002$, $\beta_1 = 0.9$ and $\beta_2 = 0.999$. With
$\beta_1^t$ we denote $\beta_1$ to the power $t$. Here,
$(\alpha/(1 - \beta_1^t))$ is the learning rate with the
bias-correction term for the first moment. All operations on
vectors are element-wise.\\[0.6em]
\textbf{Require:} $\alpha$: Stepsize\\
\textbf{Require:} $\beta_1, \beta_2 \in [0, 1)$: Exponential decay rates\\
\textbf{Require:} $f(\theta)$: Stochastic objective function with parameters $\theta$\\
\textbf{Require:} $\theta_0$: Initial parameter vector\\[0.2em]
\hspace*{1em}$m_0 \leftarrow 0$ \hfill (Initialize $1^{\text{st}}$ moment vector)\\
\hspace*{1em}$u_0 \leftarrow 0$ \hfill (Initialize the exponentially weighted infinity norm)\\
\hspace*{1em}$t \leftarrow 0$ \hfill (Initialize timestep)\\
\hspace*{1em}\textbf{while} $\theta_t$ not converged \textbf{do}\\
\hspace*{2em}$t \leftarrow t + 1$\\
\hspace*{2em}$g_t \leftarrow \nabla_\theta f_t(\theta_{t-1})$ \hfill (Get gradients w.r.t.\ stochastic objective at timestep $t$)\\
\hspace*{2em}$m_t \leftarrow \beta_1 \cdot m_{t-1} + (1 - \beta_1) \cdot g_t$ \hfill (Update biased first moment estimate)\\
\hspace*{2em}$u_t \leftarrow \max(\beta_2 \cdot u_{t-1}, |g_t|)$ \hfill (Update the exponentially weighted infinity norm)\\
\hspace*{2em}$\theta_t \leftarrow \theta_{t-1} - (\alpha/(1 - \beta_1^t)) \cdot m_t / u_t$ \hfill (Update parameters)\\
\hspace*{1em}\textbf{end while}\\
\hspace*{1em}\textbf{return} $\theta_t$ (Resulting parameters)
}}

\vspace{1em}

% ---------- Derivation ----------
\noindent
Note that the decay term is here equivalently parameterised as
$\beta_2^p$ instead of $\beta_2$. Now let $p \to \infty$, and
define $u_t = \lim_{p \to \infty}(v_t)^{1/p}$, then:
%
\begin{align}
u_t &= \lim_{p \to \infty}(v_t)^{1/p}
     = \lim_{p \to \infty}\!\left((1 - \beta_2^p) \sum_{i=1}^{t} \beta_2^{p(t-i)} \cdot |g_i|^p\right)^{\!1/p} \\
    &= \lim_{p \to \infty}(1 - \beta_2^p)^{1/p}
       \left(\sum_{i=1}^{t} \beta_2^{p(t-i)} \cdot |g_i|^p\right)^{\!1/p} \\
    &= \lim_{p \to \infty}\!\left(\sum_{i=1}^{t} \left(\beta_2^{(t-i)} \cdot |g_i|\right)^{p}\right)^{\!1/p} \\
    &= \max\!\left(\beta_2^{t-1}|g_1|,\, \beta_2^{t-2}|g_2|,\, \ldots,\, \beta_2 |g_{t-1}|,\, |g_t|\right)
\end{align}
%
Which corresponds to the remarkably simple recursive formula:
%
\begin{equation}
u_t = \max(\beta_2 \cdot u_{t-1},\, |g_t|)
\end{equation}
%
with initial value $u_0 = 0$. Note that, conveniently enough, we
don't need to correct for initialization bias in this case. Also
note that the magnitude of parameter updates has a simpler bound
with AdaMax than Adam, namely: $|\Delta_t| \leq \alpha$.

\vspace{0.5em}

\subsection*{7.2 \textsc{Temporal averaging}}

Since the last iterate is noisy due to stochastic approximation,
better generalization performance is often achieved by
averaging. Previously in Moulines \& Bach (2011), Polyak-Ruppert
averaging (Polyak \& Juditsky, 1992; Ruppert, 1988) has been
shown to improve the convergence of standard SGD, where
$\bar{\theta}_t = \frac{1}{t} \sum_{k=1}^{n} \theta_k$.
Alternatively, an exponential moving average over the parameters
can be used, giving higher weight to more recent parameter
values. This can be trivially implemented by adding one line to
the inner loop of algorithms~1 and~2:
$\bar{\theta}_t \leftarrow \beta_2 \cdot \bar{\theta}_{t-1} + (1 - \beta_2) \theta_t$,
with $\bar{\theta}_0 = 0$. Initalization bias can again be
corrected by the estimator
$\widehat{\theta}_t = \bar{\theta}_t / (1 - \beta_2^t)$.

\vspace{0.5em}

\section*{8 \textsc{Conclusion}}

We have introduced a simple and computationally efficient
algorithm for gradient-based optimization of stochastic objective
functions. Our method is aimed towards machine learning problems
with

% [page ends here — continues on page 10]

\newpage

% =========== Page 10: test_paper_page10.tex ===========
\setcounter{equation}{12}

% ---------- End of Conclusion ----------
\noindent
large datasets and/or high-dimensional parameter spaces. The
method combines the advantages of two recently popular
optimization methods: the ability of AdaGrad to deal with sparse
gradients, and the ability of RMSProp to deal with non-stationary
objectives. The method is straightforward to implement and
requires little memory. The experiments confirm the analysis on
the rate of convergence in convex problems. Overall, we found
Adam to be robust and well-suited to a wide range of non-convex
optimization problems in the field machine learning.

\vspace{0.5em}

\section*{9 \textsc{Acknowledgments}}

This paper would probably not have existed without the support
of Google Deepmind. We would like to give special thanks to Ivo
Danihelka, and Tom Schaul for coining the name Adam. Thanks to
Kai Fan from Duke University for spotting an error in the
original AdaMax derivation. Experiments in this work were partly
carried out on the Dutch national e-infrastructure with the
support of SURF Foundation. Diederik Kingma is supported by the
Google European Doctorate Fellowship in Deep Learning.

\vspace{0.5em}

\section*{\textsc{References}}

\begin{list}{}{%
  \setlength{\leftmargin}{1.5em}%
  \setlength{\itemindent}{-1.5em}%
  \setlength{\itemsep}{0.5em}%
  \setlength{\parsep}{0pt}%
}

\item Amari, Shun-Ichi. Natural gradient works efficiently in
learning. \emph{Neural computation}, 10(2):251--276, 1998.

\item Deng, Li, Li, Jinyu, Huang, Jui-Ting, Yao, Kaisheng, Yu,
Dong, Seide, Frank, Seltzer, Michael, Zweig, Geoff, He, Xiaodong,
Williams, Jason, et al. Recent advances in deep learning for
speech research at microsoft. \emph{ICASSP 2013}, 2013.

\item Duchi, John, Hazan, Elad, and Singer, Yoram. Adaptive
subgradient methods for online learning and stochastic
optimization. \emph{The Journal of Machine Learning Research},
12:2121--2159, 2011.

\item Graves, Alex. Generating sequences with recurrent neural
networks. \emph{arXiv preprint arXiv:1308.0850}, 2013.

\item Graves, Alex, Mohamed, Abdel-rahman, and Hinton, Geoffrey.
Speech recognition with deep recurrent neural networks. In
\emph{Acoustics, Speech and Signal Processing (ICASSP), 2013
IEEE International Conference on}, pp.~6645--6649. IEEE, 2013.

\item Hinton, G.E. and Salakhutdinov, R.R. Reducing the
dimensionality of data with neural networks. \emph{Science},
313(5786):504--507, 2006.

\item Hinton, Geoffrey, Deng, Li, Yu, Dong, Dahl, George E,
Mohamed, Abdel-rahman, Jaitly, Navdeep, Senior, Andrew,
Vanhoucke, Vincent, Nguyen, Patrick, Sainath, Tara N, et al.
Deep neural networks for acoustic modeling in speech
recognition: The shared views of four research groups.
\emph{Signal Processing Magazine, IEEE}, 29(6):82--97, 2012a.

\item Hinton, Geoffrey E, Srivastava, Nitish, Krizhevsky, Alex,
Sutskever, Ilya, and Salakhutdinov, Ruslan R. Improving neural
networks by preventing co-adaptation of feature detectors.
\emph{arXiv preprint arXiv:1207.0580}, 2012b.

\item Kingma, Diederik P and Welling, Max. Auto-Encoding
Variational Bayes. In \emph{The 2nd International Conference on
Learning Representations (ICLR)}, 2013.

\item Krizhevsky, Alex, Sutskever, Ilya, and Hinton, Geoffrey E.
Imagenet classification with deep convolutional neural networks.
In \emph{Advances in neural information processing systems},
pp.~1097--1105, 2012.

\item Maas, Andrew L, Daly, Raymond E, Pham, Peter T, Huang,
Dan, Ng, Andrew Y, and Potts, Christopher. Learning word vectors
for sentiment analysis. In \emph{Proceedings of the 49th Annual
Meeting of the Association for Computational Linguistics: Human
Language Technologies-Volume 1}, pp.~142--150. Association for
Computational Linguistics, 2011.

\item Moulines, Eric and Bach, Francis R. Non-asymptotic
analysis of stochastic approximation algorithms for machine
learning. In \emph{Advances in Neural Information Processing
Systems}, pp.~451--459, 2011.

\item Pascanu, Razvan and Bengio, Yoshua. Revisiting natural
gradient for deep networks. \emph{arXiv preprint
arXiv:1301.3584}, 2013.

\item Polyak, Boris T and Juditsky, Anatoli B. Acceleration of
stochastic approximation by averaging. \emph{SIAM Journal on
Control and Optimization}, 30(4):838--855, 1992.

\end{list}

% [page ends here — continues on page 11]

\newpage

% =========== Page 11: test_paper_page11.tex ===========
\setcounter{equation}{12}

\begin{list}{}{%
  \setlength{\leftmargin}{1.5em}%
  \setlength{\itemindent}{-1.5em}%
  \setlength{\itemsep}{0.5em}%
  \setlength{\parsep}{0pt}%
}

\item Roux, Nicolas L and Fitzgibbon, Andrew W. A fast natural
newton method. In \emph{Proceedings of the 27th International
Conference on Machine Learning (ICML-10)}, pp.~623--630, 2010.

\item Ruppert, David. Efficient estimations from a slowly
convergent robbins-monro process. Technical report, Cornell
University Operations Research and Industrial Engineering, 1988.

\item Schaul, Tom, Zhang, Sixin, and LeCun, Yann. No more pesky
learning rates. \emph{arXiv preprint arXiv:1206.1106}, 2012.

\item Sohl-Dickstein, Jascha, Poole, Ben, and Ganguli, Surya.
Fast large-scale optimization by unifying stochastic gradient
and quasi-newton methods. In \emph{Proceedings of the 31st
International Conference on Machine Learning (ICML-14)},
pp.~604--612, 2014.

\item Sutskever, Ilya, Martens, James, Dahl, George, and Hinton,
Geoffrey. On the importance of initialization and momentum in
deep learning. In \emph{Proceedings of the 30th International
Conference on Machine Learning (ICML-13)}, pp.~1139--1147, 2013.

\item Tieleman, T. and Hinton, G. Lecture 6.5 -- RMSProp,
COURSERA: Neural Networks for Machine Learning. Technical
report, 2012.

\item Wang, Sida and Manning, Christopher. Fast dropout training.
In \emph{Proceedings of the 30th International Conference on
Machine Learning (ICML-13)}, pp.~118--126, 2013.

\item Zeiler, Matthew D. Adadelta: An adaptive learning rate
method. \emph{arXiv preprint arXiv:1212.5701}, 2012.

\item Zinkevich, Martin. Online convex programming and
generalized infinitesimal gradient ascent. 2003.

\end{list}

% [page ends here — continues on page 12]

\newpage

% =========== Page 12: test_paper_page12.tex ===========
\setcounter{equation}{12}

\section*{10 \textsc{Appendix}}

\subsection*{10.1 \textsc{Convergence proof}}

\noindent\textbf{Definition 10.1.} \emph{A function
$f : \mathbb{R}^d \to \mathbb{R}$ is convex if for all
$x, y \in \mathbb{R}^d$, for all $\lambda \in [0, 1]$,}
\[
\lambda f(x) + (1 - \lambda) f(y) \geq f(\lambda x + (1 - \lambda) y)
\]

\noindent
Also, notice that a convex function can be lower bounded by a
hyperplane at its tangent.

\noindent\textbf{Lemma 10.2.} \emph{If a function
$f : \mathbb{R}^d \to \mathbb{R}$ is convex, then for all
$x, y \in \mathbb{R}^d$,}
\[
f(y) \geq f(x) + \nabla f(x)^\mathsf{T}(y - x)
\]

\noindent
The above lemma can be used to upper bound the regret and our
proof for the main theorem is constructed by substituting the
hyperplane with the Adam update rules.

The following two lemmas are used to support our main theorem.
We also use some definitions simplify our notation, where
$g_t \triangleq \nabla f_t(\theta_t)$ and $g_{t,i}$ as the
$i^\text{th}$ element. We define $g_{1:t,i} \in \mathbb{R}^t$ as
a vector that contains the $i^\text{th}$ dimension of the
gradients over all iterations till $t$,
$g_{1:t,i} = [g_{1,i}, g_{2,i}, \ldots, g_{t,i}]$.

\noindent\textbf{Lemma 10.3.} \emph{Let $g_t = \nabla f_t(\theta_t)$
and $g_{1:t}$ be defined as above and bounded,
$\|g_t\|_2 \leq G$, $\|g_t\|_\infty \leq G_\infty$. Then,}
\[
\sum_{t=1}^{T} \sqrt{\frac{g_{t,i}^2}{t}} \leq 2 G_\infty \|g_{1:T,i}\|_2
\]

\noindent\emph{Proof.} We will prove the inequality using
induction over $T$.

The base case for $T = 1$, we have
$\sqrt{g_{1,i}^2} \leq 2 G_\infty \|g_{1,i}\|_2$.

For the inductive step,
\begin{align*}
\sum_{t=1}^{T} \sqrt{\frac{g_{t,i}^2}{t}}
 &= \sum_{t=1}^{T-1} \sqrt{\frac{g_{t,i}^2}{t}} + \sqrt{\frac{g_{T,i}^2}{T}} \\
 &\leq 2 G_\infty \|g_{1:T-1,i}\|_2 + \sqrt{\frac{g_{T,i}^2}{T}} \\
 &= 2 G_\infty \sqrt{\|g_{1:T,i}\|_2^2 - g_T^2} + \sqrt{\frac{g_{T,i}^2}{T}}
\end{align*}

\noindent
From,
$\|g_{1:T,i}\|_2^2 - g_{T,i}^2 + \dfrac{g_{T,i}^4}{4\|g_{1:T,i}\|_2^2}
 \geq \|g_{1:T,i}\|_2^2 - g_{T,i}^2$,
we can take square root of both side and have,
\begin{align*}
\sqrt{\|g_{1:T,i}\|_2^2 - g_{T,i}^2}
 &\leq \|g_{1:T,i}\|_2 - \frac{g_{T,i}^2}{2 \|g_{1:T,i}\|_2} \\
 &\leq \|g_{1:T,i}\|_2 - \frac{g_{T,i}^2}{2 \sqrt{T G_\infty^2}}
\end{align*}

\noindent
Rearrange the inequality and substitute the
$\sqrt{\|g_{1:T,i}\|_2^2 - g_{T,i}^2}$ term,
\[
G_\infty \sqrt{\|g_{1:T,i}\|_2^2 - g_T^2}
 + \sqrt{\frac{g_{T,i}^2}{T}}
 \leq 2 G_\infty \|g_{1:T,i}\|_2
\]
\hfill$\square$

% [page ends here — continues on page 13]

\newpage

% =========== Page 13: test_paper_page13.tex ===========
\setcounter{equation}{12}

\noindent\textbf{Lemma 10.4.} \emph{Let
$\gamma \triangleq \dfrac{\beta_1^2}{\sqrt{\beta_2}}$. For
$\beta_1, \beta_2 \in [0, 1)$ that satisfy
$\dfrac{\beta_1^2}{\sqrt{\beta_2}} < 1$ and bounded $g_t$,
$\|g_t\|_2 \leq G$, $\|g_t\|_\infty \leq G_\infty$, the following
inequality holds}
\[
\sum_{t=1}^{T} \frac{\widehat{m}_{t,i}^2}{\sqrt{t\,\widehat{v}_{t,i}}}
 \leq \frac{2}{1 - \gamma}\,\frac{1}{\sqrt{1 - \beta_2}}\,\|g_{1:T,i}\|_2
\]

\noindent\emph{Proof.} Under the assumption,
$\dfrac{\sqrt{1 - \beta_2^t}}{(1 - \beta_1^t)^2}
 \leq \dfrac{1}{(1 - \beta_1)^2}$.
We can expand the last term in the summation using the update
rules in Algorithm~1,
\begin{align*}
\sum_{t=1}^{T} \frac{\widehat{m}_{t,i}^2}{\sqrt{t\,\widehat{v}_{t,i}}}
 &= \sum_{t=1}^{T-1} \frac{\widehat{m}_{t,i}^2}{\sqrt{t\,\widehat{v}_{t,i}}}
    + \frac{\sqrt{1 - \beta_2^T}}{(1 - \beta_1^T)^2}
      \cdot \frac{\left(\sum_{k=1}^{T}(1 - \beta_1)\beta_1^{T-k} g_{k,i}\right)^{\!2}}
                 {\sqrt{T \sum_{j=1}^{T}(1 - \beta_2)\beta_2^{T-j} g_{j,i}^2}} \\[0.4em]
 &\leq \sum_{t=1}^{T-1} \frac{\widehat{m}_{t,i}^2}{\sqrt{t\,\widehat{v}_{t,i}}}
    + \frac{\sqrt{1 - \beta_2^T}}{(1 - \beta_1^T)^2}
      \sum_{k=1}^{T} \frac{T\left((1 - \beta_1)\beta_1^{T-k} g_{k,i}\right)^{2}}
                           {\sqrt{T \sum_{j=1}^{T}(1 - \beta_2)\beta_2^{T-j} g_{j,i}^2}} \\[0.4em]
 &\leq \sum_{t=1}^{T-1} \frac{\widehat{m}_{t,i}^2}{\sqrt{t\,\widehat{v}_{t,i}}}
    + \frac{\sqrt{1 - \beta_2^T}}{(1 - \beta_1^T)^2}
      \sum_{k=1}^{T} \frac{T\left((1 - \beta_1)\beta_1^{T-k} g_{k,i}\right)^{2}}
                           {\sqrt{T(1 - \beta_2)\beta_2^{T-k} g_{k,i}^2}} \\[0.4em]
 &\leq \sum_{t=1}^{T-1} \frac{\widehat{m}_{t,i}^2}{\sqrt{t\,\widehat{v}_{t,i}}}
    + \frac{\sqrt{1 - \beta_2^T}}{(1 - \beta_1^T)^2}
      \cdot \frac{(1 - \beta_1)^2}{\sqrt{T(1 - \beta_2)}}
      \sum_{k=1}^{T} T\!\left(\frac{\beta_1^2}{\sqrt{\beta_2}}\right)^{\!T-k}
      \|g_{k,i}\|_2 \\[0.4em]
 &\leq \sum_{t=1}^{T-1} \frac{\widehat{m}_{t,i}^2}{\sqrt{t\,\widehat{v}_{t,i}}}
    + \frac{T}{\sqrt{T(1 - \beta_2)}}
      \sum_{k=1}^{T} \gamma^{T-k} \|g_{k,i}\|_2
\end{align*}

\noindent
Similarly, we can upper bound the rest of the terms in the
summation.
\begin{align*}
\sum_{t=1}^{T} \frac{\widehat{m}_{t,i}^2}{\sqrt{t\,\widehat{v}_{t,i}}}
 &\leq \sum_{t=1}^{T} \frac{\|g_{t,i}\|_2}{\sqrt{t(1 - \beta_2)}}
       \sum_{j=0}^{T-t} t\gamma^{j} \\
 &\leq \sum_{t=1}^{T} \frac{\|g_{t,i}\|_2}{\sqrt{t(1 - \beta_2)}}
       \sum_{j=0}^{T} t\gamma^{j}
\end{align*}

\noindent
For $\gamma < 1$, using the upper bound on the
arithmetic-geometric series,
$\sum_t t\gamma^{t} < \dfrac{1}{(1 - \gamma)^2}$:
\[
\sum_{t=1}^{T} \frac{\|g_{t,i}\|_2}{\sqrt{t(1 - \beta_2)}}
  \sum_{j=0}^{T} t\gamma^{j}
 \leq \frac{1}{(1 - \gamma)^2 \sqrt{1 - \beta_2}}
      \sum_{t=1}^{T} \frac{\|g_{t,i}\|_2}{\sqrt{t}}
\]

\noindent
Apply Lemma~10.3,
\[
\sum_{t=1}^{T} \frac{\widehat{m}_{t,i}^2}{\sqrt{t\,\widehat{v}_{t,i}}}
 \leq \frac{2 G_\infty}{(1 - \gamma)^2 \sqrt{1 - \beta_2}}
      \|g_{1:T,i}\|_2
\]
\hfill$\square$

\vspace{0.5em}

\noindent
To simplify the notation, we define
$\gamma \triangleq \dfrac{\beta_1^2}{\sqrt{\beta_2}}$.
Intuitively, our following theorem holds when the learning rate
$\alpha_t$ is decaying at a rate of $t^{-\frac{1}{2}}$ and first
moment running average coefficient $\beta_{1,t}$ decay
exponentially with $\lambda$, that is typically close to $1$,
e.g.\ $1 - 10^{-8}$.

\vspace{0.5em}

\noindent\textbf{Theorem 10.5.} \emph{Assume that the function
$f_t$ has bounded gradients,
$\|\nabla f_t(\theta)\|_2 \leq G$,
$\|\nabla f_t(\theta)\|_\infty \leq G_\infty$ for all
$\theta \in \mathbb{R}^d$ and distance between any $\theta_t$
generated by Adam is bounded,
$\|\theta_n - \theta_m\|_2 \leq D$,}

% [page ends here — continues on page 14]

\newpage

% =========== Page 14: test_paper_page14.tex ===========
\setcounter{equation}{12}

\noindent
\emph{$\|\theta_m - \theta_n\|_\infty \leq D_\infty$ for any
$m, n \in \{1, \ldots, T\}$, and $\beta_1, \beta_2 \in [0, 1)$
satisfy $\dfrac{\beta_1^2}{\sqrt{\beta_2}} < 1$. Let
$\alpha_t = \dfrac{\alpha}{\sqrt{t}}$ and
$\beta_{1,t} = \beta_1 \lambda^{t-1}$, $\lambda \in (0, 1)$.
Adam achieves the following guarantee, for all $T \geq 1$.}

\[
R(T) \leq \frac{D^2}{2\alpha(1 - \beta_1)}
          \sum_{i=1}^{d} \sqrt{T\,\widehat{v}_{T,i}}
        + \frac{\alpha(1 + \beta_1) G_\infty}
               {(1 - \beta_1)\sqrt{1 - \beta_2}\,(1 - \gamma)^2}
          \sum_{i=1}^{d} \|g_{1:T,i}\|_2
        + \sum_{i=1}^{d}
          \frac{D_\infty^2\, G_\infty \sqrt{1 - \beta_2}}
               {2\alpha(1 - \beta_1)(1 - \lambda)^2}
\]

\noindent\emph{Proof.} Using Lemma~10.2, we have,
\[
f_t(\theta_t) - f_t(\theta^*)
 \leq g_t^\mathsf{T}(\theta_t - \theta^*)
 = \sum_{i=1}^{d} g_{t,i}(\theta_{t,i} - \theta^*_{,i})
\]

\noindent
From the update rules presented in algorithm~1,
\begin{align*}
\theta_{t+1} &= \theta_t - \alpha_t \widehat{m}_t / \sqrt{\widehat{v}_t} \\
             &= \theta_t - \frac{\alpha_t}{1 - \beta_1^t}
                \left(\frac{\beta_{1,t}}{\sqrt{\widehat{v}_t}} m_{t-1}
                     + \frac{(1 - \beta_{1,t})}{\sqrt{\widehat{v}_t}} g_t\right)
\end{align*}

\noindent
We focus on the $i^\text{th}$ dimension of the parameter vector
$\theta_t \in \mathbb{R}^d$. Subtract the scalar $\theta^*_{,i}$
and square both sides of the above update rule, we have,
\[
(\theta_{t+1,i} - \theta^*_{,i})^2
 = (\theta_{t,i} - \theta^*_{,i})^2
 - \frac{2\alpha_t}{1 - \beta_1^t}\!\left(\frac{\beta_{1,t}}{\sqrt{\widehat{v}_{t,i}}} m_{t-1,i}
   + \left(1 - \frac{\beta_{1,t}}{\sqrt{\widehat{v}_{t,i}}}\right) g_{t,i}\right)\!(\theta_{t,i} - \theta^*_{,i})
 + \alpha_t^2\!\left(\frac{\widehat{m}_{t,i}}{\sqrt{\widehat{v}_{t,i}}}\right)^{\!2}
\]

\noindent
We can rearrange the above equation and use Young's inequality,
$ab \leq a^2/2 + b^2/2$. Also, it can be shown that
$\sqrt{\widehat{v}_{t,i}} = \sqrt{\sum_{j=1}^{t}(1-\beta_2)\beta_2^{t-j} g_{j,i}^2} / \sqrt{1 - \beta_2^t}
 \leq \|g_{1:t,i}\|_2$
and $\beta_{1,t} \leq \beta_1$. Then
\begin{align*}
g_{t,i}(\theta_{t,i} - \theta^*_{,i})
 &= \frac{(1 - \beta_1^t)\sqrt{\widehat{v}_{t,i}}}{2\alpha_t(1 - \beta_{1,t})}
    \left((\theta_{t,i} - \theta^*_{,i})^2 - (\theta_{t+1,i} - \theta^*_{,i})^2\right) \\
 &\quad + \frac{\beta_{1,t}}{(1 - \beta_{1,t})}
          \frac{\widehat{v}_{t-1,i}^{\,1/4}}{\sqrt{\alpha_{t-1}}}
          (\theta^*_{,i} - \theta_{t,i})
          \sqrt{\alpha_{t-1}}\,\frac{m_{t-1,i}}{\widehat{v}_{t-1,i}^{\,1/4}}
          + \frac{\alpha_t(1 - \beta_1^t)\sqrt{\widehat{v}_{t,i}}}{2(1 - \beta_{1,t})}
            \left(\frac{\widehat{m}_{t,i}}{\sqrt{\widehat{v}_{t,i}}}\right)^{\!2} \\
 &\leq \frac{1}{2\alpha_t(1 - \beta_1)}
       \left((\theta_{t,i} - \theta^*_{,i})^2 - (\theta_{t+1,i} - \theta^*_{,i})^2\right)\sqrt{\widehat{v}_{t,i}}
     + \frac{\beta_1}{2\alpha_{t-1}(1 - \beta_1)}(\theta^*_{,i} - \theta_{t,i})^2\sqrt{\widehat{v}_{t-1,i}} \\
 &\quad + \frac{\beta_1\alpha_{t-1}}{2(1 - \beta_1)}
          \frac{m_{t-1,i}^2}{\sqrt{\widehat{v}_{t-1,i}}}
        + \frac{\alpha_t}{2(1 - \beta_1)}
          \frac{\widehat{m}_{t,i}^2}{\sqrt{\widehat{v}_{t,i}}}
\end{align*}

\noindent
We apply Lemma~10.4 to the above inequality and derive the regret
bound by summing across all the dimensions for $i \in 1, \ldots, d$
in the upper bound of $f_t(\theta_t) - f_t(\theta^*)$ and the
sequence of convex functions for $t \in 1, \ldots, T$:
\begin{align*}
R(T)
 &\leq \sum_{i=1}^{d} \frac{1}{2\alpha_1(1 - \beta_1)}
       (\theta_{1,i} - \theta^*_{,i})^2 \sqrt{\widehat{v}_{1,i}}
     + \sum_{i=1}^{d} \sum_{t=2}^{T} \frac{1}{2(1 - \beta_1)}
       (\theta_{t,i} - \theta^*_{,i})^2
       \left(\frac{\sqrt{\widehat{v}_{t,i}}}{\alpha_t} - \frac{\sqrt{\widehat{v}_{t-1,i}}}{\alpha_{t-1}}\right) \\
 &\quad + \frac{\beta_1 \alpha G_\infty}{(1 - \beta_1)\sqrt{1 - \beta_2}\,(1 - \gamma)^2}
          \sum_{i=1}^{d} \|g_{1:T,i}\|_2
        + \frac{\alpha G_\infty}{(1 - \beta_1)\sqrt{1 - \beta_2}\,(1 - \gamma)^2}
          \sum_{i=1}^{d} \|g_{1:T,i}\|_2 \\
 &\quad + \sum_{i=1}^{d} \sum_{t=1}^{T} \frac{\beta_{1,t}}{2\alpha_t(1 - \beta_{1,t})}
          (\theta^*_{,i} - \theta_{t,i})^2 \sqrt{\widehat{v}_{t,i}}
\end{align*}

% [page ends here — continues on page 15]

\newpage

% =========== Page 15: test_paper_page15.tex ===========
\setcounter{equation}{12}

\noindent
From the assumption, $\|\theta_t - \theta^*\|_2 \leq D$,
$\|\theta_m - \theta_n\|_\infty \leq D_\infty$, we have:
\begin{align*}
R(T)
 &\leq \frac{D^2}{2\alpha(1 - \beta_1)}
       \sum_{i=1}^{d} \sqrt{T\,\widehat{v}_{T,i}}
     + \frac{\alpha(1 + \beta_1) G_\infty}
            {(1 - \beta_1)\sqrt{1 - \beta_2}\,(1 - \gamma)^2}
       \sum_{i=1}^{d} \|g_{1:T,i}\|_2
     + \frac{D_\infty^2}{2\alpha}
       \sum_{i=1}^{d} \sum_{t=1}^{t}
       \frac{\beta_{1,t}}{(1 - \beta_{1,t})}
       \sqrt{t\,\widehat{v}_{t,i}} \\[0.4em]
 &\leq \frac{D^2}{2\alpha(1 - \beta_1)}
       \sum_{i=1}^{d} \sqrt{T\,\widehat{v}_{T,i}}
     + \frac{\alpha(1 + \beta_1) G_\infty}
            {(1 - \beta_1)\sqrt{1 - \beta_2}\,(1 - \gamma)^2}
       \sum_{i=1}^{d} \|g_{1:T,i}\|_2 \\[0.2em]
 &\quad + \frac{D_\infty^2 G_\infty \sqrt{1 - \beta_2}}{2\alpha}
         \sum_{i=1}^{d} \sum_{t=1}^{t}
         \frac{\beta_{1,t}}{(1 - \beta_{1,t})} \sqrt{t}
\end{align*}

\noindent
We can use arithmetic geometric series upper bound for the last
term:
\begin{align*}
\sum_{t=1}^{t} \frac{\beta_{1,t}}{(1 - \beta_{1,t})} \sqrt{t}
 &\leq \sum_{t=1}^{t} \frac{1}{(1 - \beta_1)} \lambda^{t-1} \sqrt{t} \\
 &\leq \sum_{t=1}^{t} \frac{1}{(1 - \beta_1)} \lambda^{t-1}\, t \\
 &\leq \frac{1}{(1 - \beta_1)(1 - \lambda)^2}
\end{align*}

\noindent
Therefore, we have the following regret bound:
\[
R(T) \leq \frac{D^2}{2\alpha(1 - \beta_1)}
          \sum_{i=1}^{d} \sqrt{T\,\widehat{v}_{T,i}}
        + \frac{\alpha(1 + \beta_1) G_\infty}
               {(1 - \beta_1)\sqrt{1 - \beta_2}\,(1 - \gamma)^2}
          \sum_{i=1}^{d} \|g_{1:T,i}\|_2
        + \sum_{i=1}^{d}
          \frac{D_\infty^2 G_\infty \sqrt{1 - \beta_2}}
               {2\alpha \beta_1 (1 - \lambda)^2}
\]
\hfill$\square$

% [end of paper — page 15 of 15]

\end{document}
